\section{Requirements, Impacts and
Constraints}
This thesis presents a multi-modal deep learning framework for the identification of marine vessels and the tracking of vessel trajectories. It addresses three key research challenges: small vessel detection for vessels of length ¡20 m, false alarms caused by ocean clutter, and temporal consistency in trajectory reconstruction. The proposed system brings optical satellite imagery and SAR data under the same pipeline and attempts to address some fundamental limitations that cannot be effectively addressed by a single-modality approach alone. \\
Requirements, Impacts and
Constraints

In this framework, optical images are first processed by CLAHE and bilateral filtering. They are then sent through ELAN-ODConv backbone, ASFFPredict-BiFPN neck and ELANH-SPDC head, where Rotation Head and WIoU loss are used for oriented detection. On the other hand, decibel normalization and speckle reduction are applied in the SAR branch and the data is sent to a Dual-Path Feature Separation module. There, a confidence-scored bounding box is generated using physics-constrained signal backtracking and M2Det-CBAM fusion. The outputs from the two branches are then combined using IoU-based fusion and verified using Euclidean distance-based temporal consistency checking. Finally, a physics-guided DDPM trajectory imputation is performed taking into account kinematic and environmental constraints, where the trajectory of the dark vessel is reconstructed in the absence of AIS signal using dual-branch SD-KG and stochastic encoder-decoder. \\


Although full experimental validation of this framework remains a critical next step, every architectural decision in the entire pipeline is based on relevant peer-reviewed research. In addition, this overall design presents a coherent and logical solution to the limitations and research gaps that remain in previous optical, SAR, and trajectory modeling studies. 

\section{Future Work Plan} 
Taking into account the identified research gaps, the proposed methodology, and the architectural decisions outlined throughout this thesis, the future work plan is organized into four interrelated phases. The main objective of these phases is to 
transform the current framework design into a fully validated and experimentally evidence-based system. 

\subsection{Implementation and Integration of the Core Pipeline}
The first and most important task will be to implement the complete end-to-end implementation of the proposed pipeline. This will include coding of the optical preprocessing stack. In addition, the ELAN-ODConv backbone and ELANH-SPDC detection head with ASFFPredict-BiFPN neck will be trained using established remote sensing ship datasets such as HRSC2016, DOTA and SAR-Ship. At the same time, the performance of Rotation Head and WIoU loss for oriented bounding box regression will also be verified. On the other hand, the SAR branch will also be implemented at the same time. It will include an Adaptive Dual-Path Feature Separation module, an intensity-map based YOLOv8 detector, a physics-constrained signal backtracking component, and a texture map created using Root Mean Squared Deviation (RMSD). Then, an effective multi-modal ship detection system will be built by fusing two branches with an IoU-based fusion module and a Euclidean distance-based temporal consistency checker.

\subsection{Trajectory Reconstruction Validation}
After validation of the detection backbone, the next step will be on trajectory reconstruction. This phase will train and test the physics-guided DDPM on AIS datasets such as Busan Port and Norwegian coastal waters. In parallel, the ground-truth trajectory will be used as a benchmark for evaluating the accuracy of the system in filling in missing information at different levels of signal dropout. In addition, the behavioral characteristics of the vessel and the interaction rule of COLREGS will be added to the Ship-Domain Knowledge Graph. Then, through integration testing, it will be seen whether using the SD-KG output and stochastic encoder-decoder distribution together can obtain more accurate and realistic trajectories than using each branch separately. To evaluate the performance of the system, RMSE will be used to measure the position accuracy and the displacement error metric will be used to measure the effect of long-term signal absence. 

\subsection{Ablation Studies and Component Analysis}
A well-structured ablation experiment series will be conducted to determine the contribution of each architectural decision separately. A variety of experiments are planned for this purpose. These include performance comparisons of optical only, SAR-only, and fused detection to measure the improvement of multi-modal integration for different ship sizes. In addition, experiments will be conducted with and without CLAHE and the filtering preprocessing stack to evaluate their impact on the recall and false positive rate of small ship detection in high sea-state conditions. The Rotation Head will be tested with and without the orientated bounding box to evaluate whether or not the accuracy has improved over the baseline of the horizontal only bounding box. Similarly, we will test the temporal consistency module with and without it to see if it can help reduce false positives between consecutive frames. Additionally, the results of physics-guided DDPM trajectory imputation will be compared with the results of interpolation and data-driven baselines to verify the benefits of kinematic and environmental conditioning in dark vessel conditions.
 
\subsection{Real-World Deployment Considerations}
The final phase of the work will address the practical issues required to deploy a maritime surveillance system in a real-world environment.This involves the ship monitoring system’s real-time processing essentials linked with the evaluation of interference latency and computational efficiency. Additionally, to decrease the added computational overhead by the multi-modal architectures, efficiency of model compression technologies including knowledge extraction and structured shearing will be examined. Moreover, the trained model will be tested on various satellite platforms imagery that are not used in the training to evaluate cross-sensor generalization. These procedures will mirror the situation of the real-world where surveillance systems collaborate with diverse data sources. At the same time, the robustness of the model will be systematically evaluated under adverse imaging conditions such as heavy cloud cover, high sea state, and low-resolution geostationary imagery, so that potential failure modes can be identified and necessary directions for future development can be obtained. If these steps were successfully accomplished, the proposed framework will not be limited to a theoretically designed design but it will be a completely verified and test supported effective system. The main objective is to build a maritime surveillance pipeline that is resilient to a diverse set of vessel sizes and environmental situations, computationally feasible for real operational use, and capable of reconstructing the trajectory of vessels that actively try to evade detection, thereby making a meaningful contribution to the improvement of maritime safety and security. \\





